% Generated draft: requires mathematical and source review.
\documentclass[11pt]{article}
\usepackage[margin=25mm]{geometry}
\usepackage{fontspec}
\setmainfont{Latin Modern Roman}
\usepackage{amsmath,amssymb,mathtools,unicode-math}
\setmathfont{Latin Modern Math}
\usepackage{enumitem,needspace,xurl,hyperref}
\setlength{\parindent}{0pt}
\setlength{\parskip}{5pt}
\setlength{\emergencystretch}{3em}
\newcommand{\N}{\mathbb{N}}
\newcommand{\Z}{\mathbb{Z}}
\newcommand{\Q}{\mathbb{Q}}
\newcommand{\R}{\mathbb{R}}
\newcommand{\C}{\mathbb{C}}
\newcommand{\F}{\mathbb{F}}
\newcommand{\A}{\mathbb{A}}
\newcommand{\PP}{\mathbb{P}}
\newcommand{\E}{\mathbb{E}}
\newcommand{\eps}{\varepsilon}
\newcommand{\dd}{\,\mathrm{d}}
\begin{document}
\section*{UnsolvedMath ID: 30002063}
\textbf{Source problem code:} catalogue entry
\subsection{Definitions and mathematical statement}
A convex polytope is the convex hull of a finite set of points in a real Euclidean space; its dimension is the dimension of its affine hull. A simplicial polytope is a convex polytope whose proper faces are simplices. For a simplicial polytope \(P\), let \(\partial P\) denote its boundary complex, consisting of the vertex sets of its proper faces, including the empty face. Decomposability of \(P\) is interpreted below as decomposability of \(\partial P\).

A finite abstract simplicial complex \(\Delta\) is a family of subsets of a finite vertex set, containing the empty set and closed under taking subsets. Its facets are its inclusion-maximal faces. It is pure of dimension \(r\) if every facet has cardinality \(r+1\). A simplex is the complex of all subsets of a finite set; the simplex on the empty set has dimension \(-1\). For a nonempty face \(\tau\in\Delta\), define
\[
\Delta\setminus\tau=\{\sigma\in\Delta:\tau\nsubseteq\sigma\},\qquad
\operatorname{lk}_{\Delta}(\tau)=\{\sigma\in\Delta:\sigma\cap\tau=\varnothing,\ \sigma\cup\tau\in\Delta\}.
\]
Here deletion removes all faces containing \(\tau\), and \(\dim\tau=|\tau|-1\).

Use the following conventional pure-complex definitions, which are not supplied explicitly in the source record and require confirmation against the cited report. For an integer \(k\ge0\), a pure \(r\)-dimensional complex \(\Delta\) is \emph{weakly \(k\)-decomposable} if it is a simplex, or it has a nonempty face \(\tau\) with \(\dim\tau\le k\) such that \(\Delta\setminus\tau\) is pure of dimension \(r\) and weakly \(k\)-decomposable. It is \emph{\(k\)-decomposable} if it is a simplex, or it has such a face \(\tau\) for which \(\Delta\setminus\tau\) is pure of dimension \(r\) and \(k\)-decomposable, and \(\operatorname{lk}_{\Delta}(\tau)\) is pure of dimension \(r-|\tau|\) and \(k\)-decomposable. These definitions are recursive on the number of faces. To interpret the predecessor condition when \(k=0\), adopt the explicit convention that a complex is weakly \((-1)\)-decomposable, or \((-1)\)-decomposable, exactly when it is a simplex.

Let \(\mathcal P_d\) be the class of all simplicial convex polytopes of dimension \(d\), with no restriction on their number of vertices. On the usual positive-dimensional interpretation \(d\in\Z_{\ge1}\), determine separately the two thresholds
\[
\begin{aligned}
k_{\mathrm w}(d)&=\min\{k\in\Z_{\ge0}:\ \forall P\in\mathcal P_d,\ \partial P\text{ is weakly }k\text{-decomposable}\},\\
k_{\mathrm s}(d)&=\min\{k\in\Z_{\ge0}:\ \forall P\in\mathcal P_d,\ \partial P\text{ is }k\text{-decomposable}\}.
\end{aligned}
\]
The requested sharpness conditions are, respectively,
\[
\begin{aligned}
&\exists P_{\mathrm w}\in\mathcal P_d:\ \partial P_{\mathrm w}\text{ is not weakly }(k_{\mathrm w}(d)-1)\text{-decomposable},\\
&\exists P_{\mathrm s}\in\mathcal P_d:\ \partial P_{\mathrm s}\text{ is not }(k_{\mathrm s}(d)-1)\text{-decomposable}.
\end{aligned}
\]
Each minimum is understood only when its defining set is nonempty; no existence proof or value is asserted here. The source does not explicitly specify the range of \(d\), the allowed range of \(k\), the negative-index convention, or the recursive definitions. Consequently, the conventions above make the intended formulation explicit but remain subject to source review, particularly if dimension zero is intended.
\subsection{Short English statement}
For each fixed dimension \(d\), find separately the smallest decomposition parameter that works for the boundary of every simplicial convex \(d\)-polytope in the weak and ordinary senses. In each case, require a polytope for which the preceding parameter fails. The dimensional and recursive conventions stated above require confirmation against the source.
\subsection{Sources}
\begin{itemize}
\item Local UnsolvedMath registry: supplied complete record with permanent ID \(30002063\), title \emph{Decomposability Thresholds for Simplicial Polytopes}, and exact source problem code \texttt{catalogue entry}. The record belongs to the Oberwolfach Reports Open Problems collection. No individual proposer is supplied; the proposed year is recorded as 2012.
\item Supplied report citation: \emph{Triangulations (2013)}, Oberwolfach Reports, DOI \texttt{10.4171/owr/2012/24}, pp.\ 1405--1486. \url{https://doi.org/10.4171/owr/2012/24}. The background describes Oberwolfach Reports as documenting mathematical workshops and their open-problem discussions. The literature-source list also labels this reference ``original OWR report.''
\item Supplied literature reference, with its label retained as supplied: ``Hah nle--Klee--Pilaud.'' \url{https://doi.org/10.1090/S0002-9939-2014-12101-0}. No article title or full bibliographic details are supplied.
\item Supplied literature assessment dated \texttt{2026-08-21}: the record labels the problem ``partially solved'' and states that counterexamples to weak \(k\)-decomposability exist for every fixed \(k\) when dimension may grow, while the exact minimum as a function of fixed dimension \(d\) is unknown. This is supplied dated curation derived from the OWR material, not an independently verified literature assessment. The record explicitly requests independent verification of time-sensitive status and cited claims.
\end{itemize}
\end{document}
